% id: 69-01
% book: 베이직쎈 미적분1 · 04 도함수의 활용 (1)
% section: 개념 29 롤의 정리
% passage: 다음 함수 $f(x)$에 대하여 주어진 구간에서 롤의 정리를 만족시키는 상수 $c$의 값을 구하시오.
% figure: none
% answer: 
% answer_source: 
% variant_level: 0 (원본 전사)
% 이 파일은 build.mjs 가 items.json 에서 만든다. 고칠 때는 items.json 을 고친다.
\smash{\raisebox{1.5mm}{\dmkichul{베이직쎈 미적분1 69쪽 01번}}}
$f(x)=x^2-3x+8$ \qquad $[0{,}\mkern3mu\mathopen{}3]$
\dmhint{함수 $f(x)=x^2-3x+8$은 닫힌구간 $[0{,}\mkern3mu\mathopen{}3]$에서 \blank[3em]이고 열린구간 $(0{,}\mkern3mu\mathopen{}3)$에서 \blank[4em]하며 $f(0)=\blank[1.4em]=8$이므로 롤의 정리에 의하여 $f'(c)=\blank[1.2em]$인 $c$가 열린구간 $(0{,}\mkern3mu\mathopen{}3)$에 적어도 하나 존재한다.\\ 이때 $f'(x)=2x-3$이므로\\ \hspace*{1.4em}$f'(c)=2c-3=\blank[1.2em]$ \qquad $\therefore c=\blank[1.4em]$}
